\section{Related work}
\label{sec:related}

\paragraph{Societies of language-model agents.} Multi-agent language systems have
been used to simulate populations and to study what emerges from their
interaction: generative agents with memory and social behaviour
\citep{park2022social,park2023generative}, debate procedures that improve
factuality \citep{du2023debate}, opinion exchange over networks
\citep{chuang2024simulating}, and language models as stand-ins for survey
respondents \citep{argyle2023out,aher2023using}, with systematic biases documented
in simulated debate \citep{taubenfeld2024systematic}. This literature is largely
empirical: it reports what a particular population does. What it lacks is a model
in which the collective outcome can be \emph{derived} and its phase boundaries
located --- which is what we supply, along with order parameters cheap enough to
measure on any such population.

\paragraph{Bias and fairness in machine learning.} The dominant framing locates
group bias in data and in single-model decisions, as representational harm
\citep{blodgett2020language,bender2021dangers} or as allocative unfairness with
formal criteria attached
\citep{dwork2012fairness,hardt2016equality,barocas2023fairness}. Our mechanism is
orthogonal and complementary: the outcome here is produced by an irrelevant
feature interacting with the \emph{learning rule} across a group of agents, and it
is undetectable by any audit that inspects agents one at a time, since a single
agent's bias is $O(d)$ while the collective state it induces is discontinuous in
$d$. That irrelevant features degrade a classifier is old; that they can
restructure a \emph{society} of classifiers is the point here.

\paragraph{Opinion dynamics and social balance.} Models of consensus and
polarisation in interacting populations run from bounded-confidence and mixing
rules \citep{deffuant2000mixing,hegselmann2002opinion} through culture
dissemination \citep{axelrod1997dissemination} to reinforcement-driven echo
chambers \citep{baumann2020modeling}; see \citet{castellano2009statistical} for a
survey. In most of that work an agent's state is a scalar or a discrete label, so
two agents either agree or do not. Representing agents as classifiers over an
issue space \citep{caticha2011agent} makes partial agreement possible, which is
what lets us ask whether distrust follows disagreement or the reverse. Our second
diagnostic, social balance over triples
\citep{heider1946attitudes,cartwright1956structural,antal2005dynamics,marvel2011continuous},
is usually applied to sign structures that are themselves the primitive objects;
here both sign structures are generated by inference, and their balance is what
separates organised bloc formation from disorder --- the distinction that
identifies the frustrated phase, which pair correlations alone would report as
nothing happening.

\paragraph{Optimal on-line learning, and affective polarisation.} The learning
rule is the optimised on-line algorithm for the student--teacher problem, given a
Bayesian reading by \citet{opper1996online} and an entropic-dynamics formulation
by \citet{caticha2020entropic}; the identification of the inferred channel-noise
level with distrust, and its promotion to a dynamical variable, is due to
\citet{alves2016distrust}. Appendix~\ref{app:genealogy} gives the fuller
genealogy. Our contribution is not the algorithm but what it does to a population
when one class-correlated feature is added. Separately, whether groups come to
dislike each other because they disagree or the reverse is contested in political
science \citep{fiorina2008polarization,iyengar2012affect,iyengar2015fear}; we do
not adjudicate it, but identify a parameter that selects which occurs.
